\section{Experimental setup}
\label{sec:experiments-setup}

This section gathers the protocol decisions that apply to every
experiment in the chapter.

\paragraph{Hardware and software.}
Experiments are run on a single NVIDIA GPU (PyTorch \texttt{cuda}
backend) where available, falling back to CPU on machines without
one. Reproducibility is enforced through the \texttt{set\_seed}
helper of Section~\ref{subsec:method-reproducibility}, which
seeds Python, NumPy, and PyTorch (CPU and CUDA) and sets
\texttt{torch.backends.cudnn.deterministic = True}. Each run
dumps a \texttt{config.yaml} artifact that round-trips through
the loader, so any reported number can be reproduced from its
own artifact directory.

\paragraph{Repetitions.}
Every reported number is the mean (plus or minus the sample
standard deviation) across $S$ random seeds. The tightness
experiments of Section~\ref{sec:experiments-tightness} use
$S = 3$; the visual Sudoku results of
Section~\ref{sec:experiments-sudoku} use $S = 5$. The
train/validation/test splits are fixed by the benchmark
builder (Chapter~\ref{chap:datasets}) and shared across all
seeds; only the optimization seed varies.

\paragraph{Metrics.}
Both evaluation losses of
Section~\ref{subsec:loss-functions} --- the per-cell normalized
Hamming loss and the per-puzzle $0/1$ loss --- are reported on
the test split.

\paragraph{Hyperparameter selection.}
For each experiment, the regularization strength $\lambda$
(\texttt{weight\_decay}), learning rate, and schedule offset are
selected on the held-out validation split (Hamming as the
monitor metric); other hyperparameters (batch size, number of
epochs, early-stopping patience) follow the defaults of
Section~\ref{sec:method-training} unless noted otherwise.

\paragraph{Learning curves.}
When a metric is plotted as a function of a varying parameter
(e.g.\ the training-set size $N$ or the regularization strength
$\lambda$), the curve shows the mean across seeds with a shaded
band giving $\pm 1$ standard deviation.
